% Compile with XeLaTeX, pdfLaTeX, or Tectonic. No custom class or external assets.
\documentclass[11pt,letterpaper]{article}
\usepackage[margin=0.65in]{geometry}
\usepackage[T1]{fontenc}
\usepackage[utf8]{inputenc}
\usepackage{lmodern}
\usepackage{microtype}
\usepackage{enumitem}
\usepackage[hidelinks,unicode]{hyperref}
\renewcommand{\familydefault}{\sfdefault}
\hypersetup{
  pdftitle={Luis Pedro Borges - Research, Machine Learning and Data Resume},
  pdfauthor={Luís Pedro Borges},
  pdfsubject={Canadian job and internship applications; Ph.D. in progress},
  pdfkeywords={Information Retrieval, Machine Learning, NLP, Neural Search, PyTorch, Python}
}
\pagestyle{plain}
\setlength{\parindent}{0pt}
\setlength{\parskip}{3pt}
\setlength{\emergencystretch}{2em}
\setlist[itemize]{leftmargin=1.25em,itemsep=3pt,topsep=4pt,parsep=0pt,partopsep=0pt}
\urlstyle{same}
\raggedright
\newcommand{\cvsection}[1]{\vspace{9pt}{\large\bfseries #1}\par\vspace{2pt}\hrule\vspace{5pt}}
\newcommand{\entry}[3]{\textbf{#1}\hfill{\small #2}\par\nopagebreak{#3}\par\nopagebreak}
\newcommand{\publication}[3]{\textbf{#1}\par #2\par \href{#3}{Paper}\par\vspace{5pt}}

\begin{document}
{\fontsize{25}{28}\selectfont\bfseries Luís Pedro Borges}\par
\vspace{3pt}
{\large Machine Learning \textbar{} Information Retrieval \textbar{} Neural Search}\par
\vspace{3pt}
Lisbon, Portugal \textbar{} +351 935 922 727 \textbar{} \href{mailto:lpborges.careers@gmail.com}{lpborges.careers@gmail.com}\par
\href{https://luispb7.github.io/}{luispb7.github.io}
\textbar{} \href{https://github.com/LuisPB7}{GitHub: LuisPB7}
\textbar{} \href{https://www.linkedin.com/in/lu\%C3\%ADs-borges-2a3462116/}{LinkedIn}
\textbar{} \href{https://scholar.google.com/citations?user=\_BkmlLUAAAAJ}{Google Scholar}\par
\vspace{5pt}
\textbf{Canada: IEC Working Holiday approved for up to 2 years.}\par
Open work permit activation on arrival; no employer sponsorship required during the permit period.

\cvsection{Profile}
Machine learning and information retrieval researcher with a \textbf{Ph.D. in progress} at Instituto Superior Técnico and Carnegie Mellon University. Experience in neural retrieval, knowledge distillation and LLM re-ranking, including large-scale benchmark experiments and search-efficiency evaluation. Seeking Research Scientist, Data Scientist, Data Engineer, ML Engineer and Search Engineer roles and internships in Canada.

\cvsection{Technical Skills}
\textbf{Programming:} Python, C++, Java, C.\par
\textbf{Machine learning:} PyTorch, TensorFlow, Keras, NumPy, SciPy, scikit-learn.\par
\textbf{Search and NLP:} learned sparse and dense retrieval, inverted indexes, PISA, passage re-ranking, knowledge distillation, LLM re-ranking, named entity recognition, fact verification.\par
\textbf{Evaluation:} MS MARCO, TREC Deep Learning and BEIR; ranking quality, recall and retrieval latency.

\cvsection{Research Experience}
\entry{Early Stage Researcher}{Mar 2025 - Present}{INESC-ID, Lisbon, Portugal}
\begin{itemize}
  \item Research knowledge-distillation methods for small language models in information retrieval.
\end{itemize}
\vspace{4pt}
\entry{Doctoral Research - Ph.D. in progress}{Sep 2019 - Present}{Instituto Superior Técnico / INESC-ID and Carnegie Mellon University}
\begin{itemize}
  \item Co-developed \textbf{KALE}, connecting learned sparse query and document representations with efficient inverted-index search; evaluated on MS MARCO, TREC and BEIR.
  \item Conducted PyTorch distillation experiments using approximately \textbf{45 million existing synthetic query-passage judgments}, with training on a single NVIDIA A100 GPU.
  \item Evaluated \textbf{EmberTiny}, a 14.5M-parameter retriever that retained \textbf{97\% of its 110M-parameter SimLM teacher's MRR@10} on MS MARCOv1 (0.403 versus 0.415).
  \item Evaluated \textbf{SCALE + SPLADE.Lex}: \textbf{44\% lower mean index-search latency} than SPLADEv3 (60 versus 108 ms), with MRR@10 improving from 0.400 to 0.415 in MS MARCOv1 dissertation experiments.
  \item Co-developed retrieval and zero-shot LLM re-ranking for Tip-of-the-Tongue search across domains; published at SIGIR 2024, with \href{https://github.com/LuisPB7/TipTongue}{public code and data}.
\end{itemize}

\newpage
{\large\bfseries Luís Pedro Borges}\hfill{\small Machine Learning \& Information Retrieval}\par

\cvsection{Additional Experience}
\entry{Teaching Assistant - Search Engines}{Jan 2022 - May 2022}{Carnegie Mellon University, Pittsburgh, United States}
\begin{itemize}
  \item Supported the Search Engines course under Professor Jamie Callan.
\end{itemize}
\vspace{4pt}
\entry{Junior Researcher}{Mar 2019 - Sep 2019; Jul 2018 - Dec 2018}{INESC-ID, Lisbon, Portugal}
\begin{itemize}
  \item Worked on string-matching tasks in natural language processing across two research appointments.
\end{itemize}
\vspace{4pt}
\entry{Visiting Junior Researcher}{Jan 2019 - Apr 2019}{National Institute of Informatics, Tokyo, Japan}
\begin{itemize}
  \item Worked on the FEVER fact-verification task.
\end{itemize}
\vspace{4pt}
\entry{Junior Researcher Intern}{Jul 2018 - Sep 2018}{Priberam, Lisbon, Portugal}
\begin{itemize}
  \item Worked on the OntoNotes 5.0 named entity recognition task.
\end{itemize}

\cvsection{Education}
\entry{Ph.D. in Language and Information Technologies - in progress}{2019 - Present}{Instituto Superior Técnico and Carnegie Mellon University}
Dual-degree programme. Advisors: Bruno Martins (IST) and Jamie Callan (CMU).\par
Dissertation research: \emph{Latent Vocabularies for Learned Sparse Retrieval}.\par
\vspace{6pt}
\entry{M.Sc. in Information Systems and Computer Engineering}{2016 - 2018}{Instituto Superior Técnico, University of Lisbon}
Thesis: \emph{Detecting Fake News Using Deep Neural Networks}.\par
\vspace{6pt}
\entry{B.Sc. in Information Systems and Computer Engineering}{2013 - 2016}{Instituto Superior Técnico, University of Lisbon}

\cvsection{Selected Publications}
\publication{Generalizable Tip-of-the-Tongue Retrieval with LLM Re-ranking}
{L. Borges, R. Jha, J. Callan and B. Martins. \emph{SIGIR}, 2024.}
{https://doi.org/10.1145/3626772.3657917}
\publication{KALE: Using a K-Sparse Projector for Lexical Expansion}
{L. Borges, B. Martins and J. Callan. \emph{ICTIR}, 2023.}
{https://doi.org/10.1145/3578337.3605131}
\publication{Assessing the Benefits of Model Ensembles in Neural Re-ranking for Passage Retrieval}
{L. Borges, B. Martins and J. Callan. \emph{ECIR}, 2021.}
{https://doi.org/10.1007/978-3-030-72240-1_19}
\href{https://luispb7.github.io/publications/}{Full publication list and open papers}
\end{document}
